\documentclass[compress]{beamer}
\usepackage[latin1]{inputenc}
\usepackage{cancel}
\usepackage{alltt}

\definecolor{ttblue}{RGB}{48,48,170}
\newdimen\topcrop
\topcrop=10cm  %alternatively 8cm if the pdf inclusions are in letter format
\newdimen\topcropBezier
\topcropBezier=19cm %alternatively 16cm if the inclusions are in letter format

\setbeamertemplate{footline}[frame number]
\title{Improving the tactics Ring and Field\\
\small With an oracle and foreign function calls}
\author{Yves Bertot\\
joint work with Thomas Portet}
\date{July 2026}
\mode<presentation>
\begin{document}

\maketitle
\begin{frame}
\frametitle{Plan}
\begin{itemize}
\item LiberAbaci objectives: easy encoding of easy math
\item Situations where the reflexive tactics are unsatisfactory
\item A solution for {\tt ring} using {\tt ring\_simplify}
\item A remaining problem with {\tt field\_simplify}
\item More computation: greatest common divisor, Gr�bner bases
\item integrating in the reflexive framework
\end{itemize}
\end{frame}
\begin{frame}
\frametitle{LiberAbaci: Using Rocq for early Math Education}
\begin{itemize}
\item Target: Mathematics, first year of higher education
\item fine-tune the ambient theory
\begin{itemize}
\item Single type of numbers, use of subsets
\item Import theorems from other theories (prospective, using Trocq)
\end{itemize}
\item Declarative progression of proofs
\begin{itemize}
\item {\tt assert} or {\tt enough}
\item {\tt calc} as inspired by {\sc Lean}
\end{itemize}
\item Rely on a well-chosen {\tt easy} tactic for justifications
\begin{itemize}
\item Solve obvious steps, but avoid being operational ({\tt rewrite})
\item {\tt ring}, {\tt field}, {\tt nra}
\end{itemize}
\end{itemize}
\end{frame}
\begin{frame}
\frametitle{Easy encoding of easy math}
\begin{itemize}
\item Declarative proof scripts as in Misar, Isar, verbose-lean
\item {\tt calc} as inspired by Lean
\item Beware of strong tactics (students don't learn)
\item Beware of weak tactics (students learn the wrong skill)
\item Focus on {\tt ring} and {\tt field}
\end{itemize}
\end{frame}
\begin{frame}
\frametitle{The context}
\begin{itemize}
\item Given the goal \(\frac \pi {\pi ^ 2 + \pi ^ 2} = \frac 1{2 \pi}\)
\begin{itemize}
  \item {\tt field} solves it and requires the proofs of
   \(\pi ^ 2 + \pi ^ 2 \neq 0\) and \(\pi \neq 0\)
\end{itemize}
\item Given the goal \(e^{\left(\frac \pi {\pi ^ 2 + \pi ^ 2}\right)} = e^{\left(\frac 1{2 \pi}\right)}\)
\begin{itemize}
\item {\tt field} fails
\item \(x \mapsto e ^ x\) is a function outside the field fragment
\end{itemize}
\item {\tt ring} has the same problem (for formulas without divisions)
\begin{itemize}
\item \(\pi ^ 2 + \pi \times \pi = 2 \pi ^ 2\) is solved
\item \(e ^ {\left(\pi ^ 2 + \pi \times \pi\right)} = e ^ {2 \pi ^ 2}\) is not
\end{itemize}
\end{itemize}
\end{frame}
\begin{frame}
\frametitle{A study of {\tt ring}}
\begin{itemize}
\item A formal language with just integer constants, \(\times\), \(+\), \(-\),
and variables is defined
\item This language is the ring of polynoms.
\item Polynoms can also be ``normalized'': ordered in variables and powers
\item Expressions outside the ring fragment considered as ``variables''
\item {\em reification}: map formulas to the formal language
\item Normalize into polynomial expressions
\item When normalized expressions are equal, conclude positively
\end{itemize}
\end{frame}
\begin{frame}
\frametitle{Examples of {\tt ring}}
\begin{itemize}
\item \(\pi \times (\pi \times 1) + \pi \times \pi\) is {\tt X * (X * 1) + X * X}\\
\begin{itemize}
\item  {\tt X} instantiated with \(\pi\)
\item Expansion of the power function
\end{itemize}
\item By contrast \(e^{\left(\pi \times (\pi \times 1) + \pi * \pi\right)} = e ^ {2 \pi ^2}\) is
{\tt X = Y},
\begin{itemize}
\item {\tt X} instanciated with \(e^{\left(\pi ^ 2 + \pi * \pi\right)}\)\\
{\tt Y} instanciated with \(e ^ {2 \pi ^2}\)
\item Two different polynomial expressions: failure
\end{itemize}
\end{itemize}
\end{frame}
\begin{frame}
\frametitle{Proposed solution: simplify arguments}
\begin{itemize}
\item loop of {\tt ring\_simplify} actions
\item until stabilization
\item finally solve by reflexivity
\end{itemize}
\end{frame}
\begin{frame}
\frametitle{Example solution for {\tt ring}}
\begin{itemize}
\item \(cos (x + cos (x + x)) = cos (cos (2 x) + (1 - (1 - x)))\)
\item collect all formulas appearing as equality members of function arguments
\begin{itemize}
\item \(cos (x + cos (x + x))\)
\item \(x + cos (x + x)\) 
\item \(x + x\)
\item \(cos (cos (2 x) + (1 - (1 - x)))\)
\item \(cos (2 x) + (1 - (1 - x))\)
\item \(2 x\)
\end{itemize}
\end{itemize}
\end{frame}
\begin{frame}
\frametitle{Normalize the expressions}
\begin{itemize}
\item Call {\tt ring\_simplify} on all these expressions
\begin{itemize}
\item \(cos (x + cos (x + x))  \rightarrow \framebox{\(cos(x + cos (x + x))\)}\)
\item \(x + cos (x + x) \rightarrow \framebox{\(x + cos (x + x)\)}\)
\item \(x + x \rightarrow \framebox{\( 2 x\)}\)
\item \(cos (cos (2 x) + (1 - (1 - x))) \rightarrow 
\framebox{\(cos (cos (2 x) + (1 - (1 - x)))\)}\)
\item \(cos (2 x) + (1 - (1 - x)) \rightarrow \framebox{\(x + cos (2 x)\)}\)
\item \(2 x \rightarrow \framebox{\(2 x\)}\)
\end{itemize}
\item Old equality : \(cos (x + cos (x + x)) = cos (cos (2 x) + (1 - (1 - x)))\)
\item New equality : \(cos (x + cos (2 x)) = cos (x + cos (2 x)))\)
\end{itemize}
\end{frame}
\begin{frame}
\frametitle{Common structure for {\tt field} and {\tt ring}}
\begin{itemize}
\item Produce ``normal forms'' and then test their equality
\item For {\tt ring}
\begin{itemize}
\item Normal forms are polynomials
\item Equality is syntactic
\end{itemize}
\item For {\tt field}
\begin{itemize}
\item Normal forms are fractions
\item Equality is cross product: \(\frac{p}{q}=\frac{r}{s} \Leftrightarrow p \times s = q \times r \wedge q \neq 0 \wedge s \neq 0\)
\item The equality test of {\tt field} relies on {\tt ring}
\end{itemize}
\end{itemize}
\end{frame}
\begin{frame}
\frametitle{Problem with {\tt field\_simplify} for our deep behavior}
\begin{itemize}
\item {\tt field\_simplify} only produces normal forms for ``fraction equality''
\item Reification relies on syntactic equality
\item \(\frac{2 \pi}{4}\) and \(\frac{1 \pi}{2}\) are equal as fractions
\item {\tt field\_simplify} does not produce the same syntactic expression
\end{itemize}
\end{frame}
\begin{frame}
\frametitle{Improve {\tt field\_simplify}}
\begin{itemize}
\item A variant that computes reduced fractions
\item Rely on external code to compute this gcd
\item Describe the infrastructure for such an extension
\end{itemize}
\end{frame}
\begin{frame}
\frametitle{Justification theorems}
\begin{itemize}
\item For {\tt field\_simplify}\\
\(\forall l~e~\psi, {\tt Fnorm}~l~e = \psi \rightarrow
{\tt PCond (condition}~\psi{\tt )} \rightarrow\)
\strut\hskip1cm\({\tt FEeval}~l~e = {\tt P}_\phi~l~({\tt norm}~({\tt num}~\psi )) /
   {\tt P}_\phi~l~({\tt norm}~ ({\tt denum}~\psi))\)
\item For our extension\\
\(\forall l~e~\psi, {\tt Fnorm}~l~e = \psi \rightarrow
{\tt PCond (condition}~\psi{\tt )} \rightarrow\)\\
\(\forall m~\nu~\delta~\gamma,
  (m\times {\tt num}~\psi = \nu \times \gamma \wedge
  m\times {\tt denum}~\phi = \delta \times \gamma) \rightarrow\) \\
\(m \neq 0 \wedge P_\phi~l~({\tt norm}~\gamma) \neq 0 \wedge 
{\tt P}_\phi~l~({\tt norm}~\delta) \neq 0
\rightarrow\)\\
\strut\hskip1cm \({\tt FEeval}~l~e = {\tt P}_\phi~l~ ({\tt norm}~\nu) /
   {\tt P}_\phi~l~({\tt norm}~ \delta)\)
\item We use an oracle to provide the values \(m\), \(\nu\), \(\delta\), and \(\gamma\)
\item The ``common divisor'' and \(m\) non-zero conditions are easy
\item The conditions of non-zero for \(\gamma\) and \(\delta\) have to be kept
\end{itemize}
\end{frame}
\begin{frame}
\frametitle{Examples of {\tt Fnorm}, {\tt norm}, {\tt FEeval}}
\begin{itemize}
\item \(1 / (\sin 1 / (\cos 1 + \cos 1))\)
\item Unknown expressions are \(\sin 1\), \(\cos 1\)
\item the ``field'' formula \(e\) is \(1 / (X_1 / (X_2 + X_2)\)
\item {\tt Fnorm}~\(e\) is \((X_2 + X_2) / X_1\) (with conditions)
\item {\tt norm} \((X_2 + X_2)\) is \(2 X_2\)
\item \({\tt P}_\phi~[\sin 1;\cos 1]~(2 X_2)\) is \(2 \cos 1\)
\end{itemize}
\end{frame}
\begin{frame}
\frametitle{Calling an oracle}
\begin{itemize}
\item An OCaml function manipulates a datatype of ring expressions
  (with variables)
\item takes as input two polynomials
\item produces as output the greatest common divisor(with rational coefficients)
\end{itemize}
\end{frame}
\begin{frame}[fragile]
\frametitle{Polynom gcd}
\begin{verbatim}
(* file gcd.mli *)
type rat = { num: int; den: int }

type poly =
  Const of rat
| Var of int
| Add of poly * poly
| Mul of poly * poly

val poly_gcd : poly -> poly -> poly
(* TODO : only one function, returning a triplet of polynomials *)
val factorize : poly -> poly -> poly * poly
\end{verbatim}
Note that this program is not protected against overflows
\end{frame}
\end{document}


%%% Local Variables: 
%%% mode: latex
%%% TeX-master: t
%%% End: 
